%-------------------------
% RESUME — Onkar Jadhav
% CANONICAL LAYOUT TEMPLATE. Frozen preamble + frozen structure.
% Tailor only inside %== TAILOR-START/END ==  slots: HEADLINE, SUMMARY,
% COURSEWORK, EXPERIENCE-BULLETS, PROJECTS, SKILLS. Everything else,
% including section order and section headings, is frozen.
% RESUME_EVIDENCE_V1: each new tailored source requires same-directory
% <stem>.matrix.md and <stem>.provenance.md sidecars. See resumes/evidence/README.md.
%------------------------

\documentclass[a4paper,11pt]{article}
\usepackage{latexsym}
\usepackage{xcolor}
\usepackage{float}
\usepackage{ragged2e}
\usepackage[empty]{fullpage}
\usepackage{tabularx}
\usepackage{titlesec}
\usepackage{geometry}
\usepackage{marvosym}
\usepackage{verbatim}
\usepackage{enumitem}
\usepackage[hidelinks]{hyperref}
\usepackage{fancyhdr}
\usepackage{multicol}
\usepackage{graphicx}
\usepackage{lmodern}
\usepackage[T1]{fontenc}
\setlength{\multicolsep}{0pt}
\pagestyle{fancy}
\fancyhf{}
\fancyfoot{}
\renewcommand{\headrulewidth}{0pt}
\renewcommand{\footrulewidth}{0pt}
\geometry{left=1.4cm, top=0.8cm, right=1.2cm, bottom=1cm}
\usepackage[most]{tcolorbox}
\tcbset{
    frame code={}
    center title,
    left=0pt, right=0pt, top=0pt, bottom=0pt,
    colback=gray!20, colframe=white,
    width=\dimexpr\textwidth\relax,
    enlarge left by=-2mm,
    boxsep=4pt, arc=0pt, outer arc=0pt,
}
\urlstyle{same}
\raggedright
\setlength{\tabcolsep}{0in}

\titleformat{\section}{
  \vspace{-4pt}\scshape\raggedright\large
}{}{0em}{}[\color{black}\titlerule \vspace{-7pt}]

\newcommand{\resumeItem}[2]{
  \item{\textbf{#1}{\hspace{0.5mm}#2 \vspace{-0.5mm}}}
}
\newcommand{\resumeSubheading}[4]{
\vspace{0.5mm}\item
    \begin{tabular*}{0.98\textwidth}[t]{l@{\extracolsep{\fill}}r}
        \textbf{#1} & \textit{\footnotesize{#4}} \\
        \textit{\footnotesize{#3}} & \footnotesize{#2}\\
    \end{tabular*}
    \vspace{-2.4mm}
}
\newcommand{\resumeProject}[4]{
\vspace{0.5mm}\item
    \begin{tabular*}{0.98\textwidth}[t]{l@{\extracolsep{\fill}}r}
        \textbf{#1} \textnormal{---} \textit{#2} & \textit{\footnotesize{#4}} \\
        \footnotesize{\textit{#3}} & \\
    \end{tabular*}
    \vspace{-2.4mm}
}

\renewcommand{\labelitemi}{$\vcenter{\hbox{\tiny$\bullet$}}$}
\newcommand{\resumeSubHeadingListStart}{\begin{itemize}[leftmargin=*,labelsep=0mm]}
\newcommand{\resumeHeadingSkillStart}{\begin{itemize}[leftmargin=*,itemsep=1.7mm, rightmargin=2ex]}
\newcommand{\resumeItemListStart}{\begin{justify}\begin{itemize}[leftmargin=3ex, rightmargin=2ex, noitemsep,labelsep=1.2mm,itemsep=0mm]\footnotesize}
\newcommand{\resumeSubHeadingListEnd}{\end{itemize}\vspace{2mm}}
\newcommand{\resumeHeadingSkillEnd}{\end{itemize}\vspace{-2mm}}
\newcommand{\resumeItemListEnd}{\end{itemize}\end{justify}\vspace{-2mm}}

\newcolumntype{L}{>{\raggedright\arraybackslash}X}
\newcolumntype{R}{>{\raggedleft\arraybackslash}X}
\newcolumntype{C}{>{\centering\arraybackslash}X}

\newcommand{\name}{Onkar Jadhav}
\newcommand{\phone}{7058604679}
\newcommand{\emaila}{onkarj012@gmail.com}
\newcommand{\github}{github.com/Onkarj012}
\newcommand{\linkedin}{linkedin.com/in/onkarj012}

\begin{document}

%----------HEADING-----------------
\begin{tabularx}{\linewidth}{L r}
  \textbf{\Large \name} & {+91-\phone} \\
%== TAILOR-START: HEADLINE ==
  {AI / Full-Stack Software Engineer} & \href{mailto:\emaila}{\emaila} \\
%== TAILOR-END: HEADLINE ==
  {Pune, Maharashtra} & \href{https://\github}{\github} \\
  \href{https://leo4.dev}{leo4.dev} & \href{https://\linkedin}{\linkedin}
\end{tabularx}
\vspace{-4mm}

%----------SUMMARY-----------------
\section{\textbf{Summary}}
\vspace{1mm}
%== TAILOR-START: SUMMARY ==
\small{
AI/full-stack engineer (B.Tech CS, PCCOE, 9.17 CGPA) building full-stack AI products end to end, from React/Next.js frontends to Python/Node backend services. Architected a hybrid RAG system (Next.js + FastAPI, BM25 + pgvector, reciprocal rank fusion) gated behind a deploy-time eval harness, and an agentic creator--critic--evaluator pipeline with cross-model benchmarking. Currently shipping Stride, a five-agent AI product, as Lead Engineer. Strong across React, Next.js, FastAPI, Node.js, RAG, and agentic AI system design.
}
%== TAILOR-END: SUMMARY ==
\vspace{-3mm}

%-----------EDUCATION-----------
\section{\textbf{Education}}
  \resumeSubHeadingListStart
    \resumeSubheading
      {Bachelor of Technology in Computer Science}{CGPA: 9.17}
      {Pimpri Chinchwad College of Engineering, Savitribai Phule Pune University}{Nov 2022 -- Jun 2026}
      \resumeItemListStart
%== TAILOR-START: COURSEWORK ==
        \item {Relevant Coursework: Machine Learning, Data Structures \& Algorithms, Database Management Systems, Operating Systems, Computer Networks, Object-Oriented Programming}
%== TAILOR-END: COURSEWORK ==
      \resumeItemListEnd
  \resumeSubHeadingListEnd
\vspace{-5.5mm}

%-----------EXPERIENCE-----------------
\section{\textbf{Experience}}
\resumeSubHeadingListStart

    \resumeSubheading
      {Stride --- AI-First Adaptive Wellness Platform}{Remote}
      {Lead Engineer}{April 2026 -- Present}
      \resumeItemListStart
%== TAILOR-START: EXPERIENCE-BULLETS ==
        \item {Architected and shipped Stride to production as sole engineer, owning the full stack --- React/Next.js frontend, backend services, and data layer --- running live with an early private user cohort driving iterative product feedback ahead of public beta.}
        \item {Architected five specialised domain agents (diet, workout, sleep, hydration, habits) sharing a behavioral memory layer, with multi-provider LLM routing via OpenRouter across Anthropic Claude, OpenAI GPT, and Gemini --- eliminating vendor lock-in across five concurrent agent types.}
        \item {Separated LLM reasoning from deterministic computation (BMR/NEAT/EAT/TEF calculator), eliminating opaque model arithmetic errors; implemented CI-style quality gates before each capability release.}
        \item {Selected Convex for real-time reactive state over a REST + WebSockets stack, keeping five concurrent agent types in sync across the frontend without manual polling logic.}
%== TAILOR-END: EXPERIENCE-BULLETS ==
      \resumeItemListEnd

\resumeSubHeadingListEnd
\vspace{-5mm}

%-----------ENGINEERING PROJECTS-----------------
\section{\textbf{Engineering Projects}}
\resumeSubHeadingListStart

%== TAILOR-START: PROJECTS ==
    \resumeProject
      {Lore Engine}{RAG Platform with Deploy-Gated Evaluation Harness}
      {Next.js, FastAPI, PostgreSQL, pgvector, BM25, Celery, Redis, Cloudflare R2, OpenAI, Anthropic Claude}
      {2026 -- Present}
      \resumeItemListStart
        \item {Architected a full-stack RAG application ingesting \textbf{6,000+ MediaWiki pages} per series into a hybrid retrieval pipeline (BM25 + pgvector via reciprocal rank fusion) --- chosen over pure vector search to handle exact entity and proper-noun matching that semantic models fail on.}
        \item {Designed section-aware chunking at H2/H3 boundaries with enriched \texttt{[Series][Page][Section]} metadata and spoiler-filter tags per chunk, demonstrating fine-grained context control over large knowledge bases.}
        \item {Established a \textbf{50-question per-series ground-truth eval harness} measuring retrieval MRR and end-to-end answer accuracy, gating every deploy --- the same CI/CD-for-AI pattern required in production LLM analytics systems.}
        \item {Built an async ingestion pipeline (Celery, Redis, Cloudflare R2) that fetches from the MediaWiki API and converts wikitext to markdown, scaling to \textbf{$\sim$5M vectors per series} in pgvector --- chosen over Pinecone/Weaviate to stay in one network hop.}
      \resumeItemListEnd
    \vspace{-2mm}

    \resumeProject
      {PromptForge}{Agentic LLM Orchestration \& Cross-Model Benchmarking Platform}
      {Node.js, Express.js, BullMQ, Convex, OpenRouter, React.js}
      {Jan 2026 -- Feb 2026}
      \resumeItemListStart
        \item {Architected a creator--critic--evaluator agentic orchestration pipeline, scoring prompt clarity, constraint coverage, and token efficiency across versioned, regression-tracked revision trails --- treating prompts as testable production artifacts.}
        \item {Implemented cross-model routing across OpenAI, Anthropic, Gemini, and open-source models via OpenRouter for objective quality, latency, and cost-per-output benchmarking.}
        \item {Built a fault-tolerant async backend using BullMQ workers with idempotent retries and Convex persistence, supporting reliable regression testing at scale without job loss on provider timeouts.}
        \item {Emitted structured prompt outputs with metadata --- token budget, latency profile, expected output shape, version changelog --- enabling prompts to be reviewed and deployed like code.}
      \resumeItemListEnd
%== TAILOR-END: PROJECTS ==

\resumeSubHeadingListEnd
\vspace{-8.5mm}

%-----------TECHNICAL SKILLS-----------------
\section{\textbf{Technical Skills}}
 \begin{itemize}[leftmargin=0.05in, label={}]
    \small{\item{
%== TAILOR-START: SKILLS ==
     \textbf{Frontend}{: React.js, Next.js, Tailwind CSS, HTML, CSS} \\
     \textbf{Backend}{: FastAPI, Node.js, Express.js, BullMQ, Celery, Redis, REST APIs, WebSockets} \\
     \textbf{AI / ML}{: LLM Orchestration, Agentic Workflows, RAG Pipelines, Deploy-Gated Evaluation, pgvector, BM25, Reciprocal Rank Fusion, Multi-Agent Systems, Prompt Engineering, Behavioral Memory, TensorFlow, OpenRouter} \\
     \textbf{Languages}{: Python, TypeScript, JavaScript, SQL, C++, Java} \\
     \textbf{Databases}{: PostgreSQL (+ pgvector), MongoDB, Convex} \\
     \textbf{Cloud \& DevOps}{: AWS, Cloudflare R2, Git, GitHub, async job pipelines, CI-style eval gating} \\
     \textbf{System Design}{: Multi-agent coordination, hybrid retrieval, event-driven architecture, snapshot-first serving} \\
%== TAILOR-END: SKILLS ==
    }}
 \end{itemize}
\vspace{-8pt}

\end{document}
